% TOP_PROBLEM: {"id": 30007132, "title": "Affine Bernstein assertion in dimensions three through nine", "category_id": 6, "category": {"id": 6, "name": "geometry", "display_name": "Geometry", "description": "Euclidean and non-Euclidean geometry, geometric structures.", "slug": "geometry", "order_index": 6, "created_at": "2026-07-31T15:26:25.671Z"}, "status": "open"}
% FIRST_POSTED: 2026-10-06
% ATTEMPT_STATUS: solved
% OPENAI_CLAIM: {"schema_version":1,"source_repo":"https://github.com/openai/math","source_commit":"adc7f1241b42e322a6451854ab7e4b4c146bf78a","release_date":"2026-10-06","match":"full","resolution":"proved","independently_reviewed":false,"second_pass":"agreed","families":[{"family":"353","title":"Affine Bernstein rigidity through dimension nine and a smooth dimension-ten counterexample","subject":"Differential geometry","manuscripts":[{"title":"The affine Bernstein theorem in dimensions three through nine","pdf_path":"preprints/The-affine-Bernstein-theorem-in-dimensions-three-through-nine-September-24-2026/main.pdf","readme_path":"preprints/The-affine-Bernstein-theorem-in-dimensions-three-through-nine-September-24-2026/README.md","date":"2026-09-24","theorem_ref":"preprints/The-affine-Bernstein-theorem-in-dimensions-three-through-nine-September-24-2026/build/main.tex#thm:main"}],"lean":{"doc_path":"lean/docs/353.md","comparators":["lean/ComparatorChallenges/AffineBernstein.json","lean/ComparatorChallenges/AffineBernstein.lean"],"declarations":[],"covers_main_theorem":true},"reasoning_trace":null}]}
\documentclass[11pt]{article}
\usepackage[margin=25mm]{geometry}
\usepackage{fontspec}
\setmainfont{Latin Modern Roman}
\usepackage{amsmath,amssymb,mathtools}
\usepackage{unicode-math}
\setmathfont{Latin Modern Math}
\usepackage{enumitem}
\usepackage{xurl,hyperref,bookmark}
\hypersetup{colorlinks=true,urlcolor=blue}
\setcounter{secnumdepth}{0}
\setlength{\parindent}{0pt}
\setlength{\parskip}{5pt}
\setlength{\emergencystretch}{3em}
\begin{document}
\section*{Affine Bernstein assertion in dimensions three through nine}

\subsection{OpenAI's claim}
\textit{Affine Bernstein rigidity through dimension nine}: The paper claims that Euclidean-complete smooth strictly convex affine-maximal graphs in dimensions three through nine are elliptic paraboloids. It also treats affine-complete immersed hypersurfaces without an initial global graph hypothesis.

\subsection{Outline of the argument}
\textit{Affine Bernstein rigidity through dimension nine}: Tangent sections are analyzed through affine limits of the epigraph. Maximizing the number of nonnegative model directions yields uniform centering of transverse fibers, since a centering failure would create a further direction after rescaling. This section balance supplies the global control needed for the affine equation’s rigidity conclusion, with a separate immersion reduction handling the nongraph setting.

\subsection{Scope relative to this problem}
\textit{Affine Bernstein rigidity through dimension nine}: The paper claims that Euclidean-complete smooth strictly convex affine-maximal graphs in dimensions three through nine are elliptic paraboloids. It also treats affine-complete immersed hypersurfaces without an initial global graph hypothesis. The parameters, regularity assumptions and normalization in this question are retained in applying that assertion.

With the hypotheses and conventions above, the claimed result addresses the whole mathematical assertion on this page.

\subsection{Formal verification}
Family 353 describes Lean coverage for the Euclidean-complete graph rigidity in dimensions three through nine; the nongraph and dimension-ten assertions are separate.

\textit{Affine Bernstein rigidity through dimension nine}: The manuscript links \texttt{AffineBernstein.lean}.

The documented selection includes this manuscript’s main assertion; that assertion retains the restrictions stated above.

This site has not run the Lean development. These links describe OpenAI’s selected formal statements and should be read with their stated scope.

\subsection{Status caveat}
OpenAI reports that an internal model produced the manuscripts. Their arguments have not been independently verified as proofs by this site. Any portion outside the documented formal scope may contain mistakes, and the catalogue label reports the announced mathematical claim.

\subsection{Sources}
\begin{itemize}
\item OpenAI, \emph{The affine Bernstein theorem in dimensions three through nine}, 2026-09-24; citation key \texttt{OAI:The-affine-Bernstein-theorem-in-dimensions-three-through-nine-September-24-2026}. \href{https://github.com/openai/math/blob/adc7f1241b42e322a6451854ab7e4b4c146bf78a/preprints/The-affine-Bernstein-theorem-in-dimensions-three-through-nine-September-24-2026/main.pdf}{PDF (GitHub)}; \href{https://github.com/openai/math/raw/adc7f1241b42e322a6451854ab7e4b4c146bf78a/preprints/The-affine-Bernstein-theorem-in-dimensions-three-through-nine-September-24-2026/main.pdf}{PDF (direct)}; \href{https://github.com/openai/math/blob/adc7f1241b42e322a6451854ab7e4b4c146bf78a/preprints/The-affine-Bernstein-theorem-in-dimensions-three-through-nine-September-24-2026/README.md}{Manuscript page and citation}; \href{https://github.com/openai/math/blob/main/preprints/The-affine-Bernstein-theorem-in-dimensions-three-through-nine-September-24-2026/main.pdf}{Latest version}.
\item Family 353: \href{https://github.com/openai/math/blob/adc7f1241b42e322a6451854ab7e4b4c146bf78a/lean/docs/353.md}{Lean scope}.
\item \href{https://github.com/openai/math/blob/adc7f1241b42e322a6451854ab7e4b4c146bf78a/lean/ComparatorChallenges/AffineBernstein.json}{Comparator AffineBernstein.json}.
\item \href{https://github.com/openai/math/blob/adc7f1241b42e322a6451854ab7e4b4c146bf78a/lean/ComparatorChallenges/AffineBernstein.lean}{Comparator AffineBernstein.lean}.
\item \href{https://github.com/openai/math/blob/adc7f1241b42e322a6451854ab7e4b4c146bf78a/CONTENTS.md}{OpenAI catalogue entry}.
\end{itemize}
\end{document}
